\section{Natural and Supernatural Happiness}

\begin{quotex}
In any case, the important thing is to make the distinction, well known to traditional teachings, between the happiness or pleasure that is ardent, and that which is heroic ... The distinction corresponds to that between two opposite attitudes and two opposite human types. The first type of happiness or pleasure belongs to the naturalistic plane and is marked by passivity toward the world of impulses, instincts, passions, and inclinations. Tradition defines the basis of naturalistic existence as desire and thirst, and ardent pleasure is that which is tied to the satisfaction of desire in terms of a momentary dampening of the fire that drives life onward. Heroic pleasure, on the other hand, is that which accompanies a decisive action that comes from Being, from the plane superior to that of life... \flright{\textsc{Julius Evola}}
\end{quotex}


\begin{quotex}
Humanity surges with uncontrolled passion, is tumultuous with ungoverned grief, is blown about by anxiety and doubt. Only the wise man, only he whose thoughts are controlled and purified, makes the winds and storms of the soul obey him.
\flright{\textsc{James Allen}, \emph{As a Man Thinketh}}
\end{quotex}


\paragraph{Natural Happiness}


Although the vast majority of people in all times and places has desired natural happiness, Tradition teaches that, in every age, there are a few who are called to transcend such desires and reach as state of peaceful bliss. Evola lists four motivating powers of desire.

\begin{description}
\item{Instinct} Inborn desires that do not need to be taught.

\item{Impulse} A strong idea that creates desire. 

\item{Passion} A strong emotion that creates desire. 

\item{Inclination} Arbitrary likes and dislikes that lead to desire or aversion.
\end{description}

Instincts, impulses, passions, and inclinations create ardent desires that the natural man seeks to quench. This is always accompanied by the belief that the satisfaction of the desire will lead to happiness. However, although the natural or carnal man may experience happiness for a season, as soon as the fire is quenched, another desire---equally ardent---arises, again demanding satisfaction.

\paragraph{Desire and Volition}


To understand the ``passive'' and ``active'' forces in the face of desire, it is necessary to observe the process of conation clearly. The \textbf{Principle of Sufficient Reason} leads us to look for the reason for things and events; for example, this may be the logical necessity of a mathematical proof or the physical cause for things in the world. In the case of human action, it is the motive, or the reason a man gives to explain what he does.

A man, then, will have a motive to reach some result and he is free to choose the means to reach that result---in this limited sense he is ``free''. However, is he free to choose the very motive that drives his actions? If a man's consciousness is passive in respect to instincts, impulses, passions, and inclinations, then his motivations are formed by those factors external to him. In this much deeper sense, such a man is not free and is in bondage to forces arising outside him.

\paragraph{Detachment}


\begin{quotex}
Detachment coexists with a fully lived experience; a calm being is constantly wedded to the substance of life. The consequence of this union, existentially speaking, is a most particular kind of lucid inebriation, one might almost say intellectualized and magnetic, which is the absolute opposite of what comes from the ecstatic opening to the world of elementary forces, instinct, and nature. In this very special inebriation, subtilized and clarified, is to be seen the vital element necessary for an existence in the free state, in a chaotic world abandoned to itself.
\flright{\textsc{J Evola}}
\end{quotex}


Supernatural happiness, or ``heroic pleasure'', or ``lucid inebriation'' is the effect of detachment. Detachment is not the same as withdrawal from life. Rather, the detached man is the active force in respect to life. No longer is he the puppet of instincts, impulses, passions, and inclinations. Instead, he is the center of awareness, the ``unmoved mover'', the observer. Instincts, impulses, passions, and inclinations become merely ``options'' and not fate, necessity, \emph{``ananke''}.


\flrightit{Posted on 2007-09-07 by Cologero}

\begin{center}* * *\end{center}

\begin{footnotesize}
\begin{sffamily}
\texttt{Me on 2016-03-23 at 15:56 said:}

``That exquisite poise of character which we call serenity is the last lesson culture; it is the flowering of life, the fruitage of the soul. It is precious as wisdom, more to be desired than gold -- yea, than even fine gold. How insignificant mere money-seeking looks in comparison with a serene life -- a life that dwells in the ocean of Truth, beneath the waves, beyond the reach of tempests, in the Eternal Calm!'' (Allen)

\url{https://www.youtube.com/watch?v=nW8M6Km\_baA}



\hfill

\end{sffamily}
\end{footnotesize}

